\section{Arsitektur Parser Generator}
\label{sec:ch13-arsitektur}

\subsection{Ide Dasar}

Parser generator menerima spesifikasi grammar dan menghasilkan kode
parser secara otomatis. Pengguna menulis aturan produksi dan
semantic action; generator membangun tabel LR atau LL dan menghasilkan
fungsi parser. Pendekatan ini memisahkan deskripsi bahasa dari
detail algoritma parsing.

\begin{figure}[H]
\centering
\begin{tikzpicture}[node distance=6mm and 9mm]
  \node[phase, text width=2.8cm] (grammar) {Grammar\\dan aksi};
  \node[phase, text width=2.8cm, right=of grammar] (gen)
        {Parser\\generator};
  \node[phase, text width=2.8cm, right=of gen] (code)
        {Kode parser\\dan tabel};
  \node[phase, text width=2.8cm, below=of code] (ast) {AST};
  \draw[arrow] (grammar) -- (gen);
  \draw[arrow] (gen) -- (code);
  \draw[arrow] (code) -- (ast);
\end{tikzpicture}
\caption{Alur kerja parser generator}
\label{fig:ch13-arsitektur}
\end{figure}

\subsection{Keuntungan}

Parser generator memberikan beberapa keuntungan:
\begin{itemize}[leftmargin=*]
  \item mengurangi kode yang ditulis tangan;
  \item tabel dibangun secara sistematis sehingga mengurangi kesalahan;
  \item perubahan grammar mudah diterapkan;
  \item mendukung grammar yang lebih besar daripada recursive descent.
\end{itemize}

Keterbatasannya adalah ketergantungan pada alat, pesan kesalahan yang
kadang kurang kontekstual, dan kurva belajar untuk membaca laporan
konflik.

\subsection{Alat yang Umum}

Alat yang dipakai pada buku ini:
\begin{itemize}[leftmargin=*]
  \item \textbf{Bison}/\textbf{Yacc}: parser generator dengan grammar
        LALR(1) dan aksi C/C++;
  \item \textbf{Flex}: lexer generator yang menghasilkan fungsi
        \code{yylex} untuk Bison;
  \item \textbf{Parser tangan}: recursive descent yang dibahas pada
        Bab X, dipakai sebagai pembanding.
\end{itemize}

Buku ini memakai Bison untuk parser dan Flex untuk lexer karena
keduanya terintegrasi baik dengan C. Parser tangan dibahas sebagai
pembanding, bukan sebagai pengganti alat utama.
